\section*{Summary}
\begin{enumerate}
    \item \textbf{Automaticity \& Core Weapons (Ch.1):} Modern complexity drives reliance on fast decision shortcuts. Persuasion often works by triggering six core rules: reciprocation, commitment/consistency, social proof, liking, authority, and scarcity.
    \item \textbf{Contrast \& Sequencing (Ch.1):} Judgments are relative; extreme anchors make later options look reasonable. Watch staged sequences (high–low pricing, feature stacks) that change perceived value without improving actual merit.
    \item \textbf{Reciprocation (Ch.2):} Uninvited gifts and \emph{reciprocal concessions} (door-in-the-face) create indebtedness. Defend by classifying the benefit, returning commensurately, naming the sequence, and pausing before agreeing.
    \item \textbf{Commitment/Consistency (Ch.3):} Public, effortful, and voluntary commitments shape self-image and future actions (e.g., foot-in-the-door). Prefer flexible self-definitions; make private, revocable commitments when stakes are high.
    \item \textbf{Social Proof (Ch.4):} Under uncertainty, we look to others—especially similar others—for cues (pluralistic ignorance, bystander effect). Defend by verifying independence, and isolating aid requests to break diffusion of responsibility.
    \item \textbf{Liking (Ch.5):} Attractiveness halos, similarity, compliments, familiarity/cooperation, and positive associations increase compliance. Separate the person from the proposal; accept only diagnostic signals tied to the decision.
    \item \textbf{Authority (Ch.6):} Titles, uniforms, and trappings trigger \emph{directed deference}. True authority is domain-specific and trustworthy. Verify credentials and prefer two-sided, transparent evidence over status signals.
    \item \textbf{Scarcity (Ch.7):} Loss aversion, deadlines, limited-number claims, reactance, competition, and secrecy heighten desire—especially when scarcity is sudden. Install time-outs; choose via pre-set criteria, not urgency.
    \item \textbf{Stacked Tactics (Ch.1–7):} Persuaders combine cues (e.g., scarcity + social proof; reciprocation + contrast; authority + association) to magnify compliance. Evaluate each cue’s \emph{diagnostic relevance} before weighing it.
    \item \textbf{Decision Safeguards (All):} Precommit standards, add friction (delays, second opinions), separate cues from merits, and re-evaluate in neutral contexts. Use ethical influence only when signals reflect genuine quality.
\end{enumerate}
